\section{The constant terms and a square-root period comparison}
\label{sec:constant-terms}

We retain the theta data of Proposition~\ref{prop:theta-family} and the
integer $b_*$ of Section~\ref{sec:tame}.  At each of the two places of $F$
above $q=q_i$, the upper theta test, in the column order $(e,B_1,B_2,f)$,
is supported on
\[
 Y_{[e,B_1]}\in\GL_2(\cO_{F_v}),\qquad
 Y_{B_2},Y_f\in\cO_{F_v}^{2},
\]
and has factor $\eta_v(\det Y_{[e,B_1]})$ on the invertible block.  Its
transformation under the full source maximal compact cancels the
determinant twist in the source integral.  The corresponding two-column
test is the reference test for the binary theta lift.  All unspecified
exceptional tests have fixed support and level and uniformly bounded
coefficients.

\begin{proposition}\label{prop:constant-term-small}
Assume that the conductor-one Heegner point $y$ is torsion.  The positive
weights $m_i$ in Proposition~\ref{prop:theta-family} can be chosen with
$m_i\to+\infty$ and $m_i\to0$ on the prescribed $p$-adic branch so that
every constant Fourier--Jacobi coefficient of the upper theta form tends
coefficientwise to zero modulo $t^{b_*}$.  The assertion holds in bounded
algebraic ordinary differential bases, uniformly over all boundary
representatives and all polynomial components.  The weights may be
required to satisfy any additional lower bounds and any further
$p$-adic approximation conditions tending to zero.
\end{proposition}

Here and below, a coefficientwise limit uses the full tame Artin
deformation of Section~\ref{sec:tame}.  In particular, it is stronger than
a statement about separate specializations at finite-order characters.

At index zero, the hyperbolic variable in the mixed expansion is zero.
Every minor involving its column vanishes, and the remaining Fock
polynomial is the determinant power on $B$.  A constant coefficient is
therefore a value of the binary theta lift.  We first identify the single
global scalar in that lift.  Its square will be expressed by two raw
toric periods, one of which approaches the Heegner logarithm from
Proposition~\ref{prop:heegner-log}.

\subsection{The binary scalar and its multiplicity}

Since multiplying a Hermitian form by the scalar $b_0\delta$ does not
change its isometry group, $U(B)=U(A)$.  The binary theta lift of the
source data belongs to the compact automorphic representation
$\pi\eta(\det)$.  We record the multiplicity assertion needed here.
At split places equal-rank local theta correspondence gives the
prescribed representation; at nonsplit finite places we use only its
spherical line, and at infinity the binary polynomial gives the
determinant character.  Removing $\eta(\det)$ gives trivial adelic
central action.

There is a direct description of the resulting automorphic space that
avoids a multiplicity formula for general inner-form packets.  Let
$D_{\mathrm{quat}}/F$ be the quaternion algebra ramified at the two real
places and split at all finite places, so that
$PU(A)=P(D_{\mathrm{quat}}^\times)$.  Write $\omega_{L/F}$ for the
quadratic idele-class character.  The image of the central quotient of
$U(A)$ in the projective quaternion automorphic quotient is the part
on which
\[
 \omega_{L/F}(\operatorname{Nrd}(g))=1.
\]
Locally the obstruction to lifting is the reduced norm modulo norms
from $L$.  The Hasse norm theorem supplies rational lifts when the local
obstructions vanish, and rational quaternion norms are precisely the
totally positive elements.  Quadratic norm reciprocity then identifies
the adelic image modulo rational translation with the displayed part.

Extend a function on this part by zero to the projective quaternion
quotient.  This respects all split-place actions and the trivial
archimedean weight.  After Jacquet--Langlands transfer, the split good
Hecke data determine the prescribed cuspidal Hilbert representation up
to the twist $\omega_{L/F}$.  To see the last assertion, use the Galois
representations of the parallel-weight-two Hilbert transfers
\cite[Theorem~1.1 and Remark~1.2]{Newton2015} and Chebotarev at the
places splitting in $L$.  Their restrictions to $G_L$
agree and are absolutely irreducible, so their extensions to $G_F$
differ by at most the quadratic character of $L/F$.  Strong multiplicity
one gives the assertion for the automorphic representations.  The
one-dimensional reduced-norm representations do not have these
cuspidal Hecke data.  The two possible quadratic twists restrict
identically to the indicated part.  Moreover restriction of the
prescribed representation is injective: a nonzero function supported
only on the complementary part would belong also to its quadratic
twist, contradicting their distinctness.  Finally each nonsplit
hyperspecial factor contributes a single spherical line.  This proves
the required multiplicity one with exactly the local levels being
used.  The base-change and Jacquet--Langlands inputs are the standard
cuspidal ones \cite{ArthurClozel1989,JacquetLanglands1970}; the avoidance
of the inducing CM field ensures the required cuspidality.

At a split finite place the binary local intertwiner is
\begin{equation}\label{eq:binary-intertwiner}
 J_v(\Phi,z)=\int_{\GL_2(F_v)}
   \Phi({}^th)|\det h|_v\,
   \pi_v^\vee(h)\eta_v^{-1}(\det h)z\,d^\times h,
 \qquad\vol(\GL_2(\cO_{F_v}))=1.
\end{equation}
Here $z$ belongs to the bare $\pi_v^\vee$ space; the determinant twist
has been written in the integrand.  The output is read by inverse
transpose, which is field conjugation in the chosen split bases.
Equation~\ref{eq:binary-intertwiner} is the Godement--Jacquet integral
at the unitary center, interpreted by rational continuation where
necessary.  Equivariance, local theta uniqueness, and a nonzero open
test identify it with the theta intertwiner.  We normalize by this
map at the exceptional split places, by spherical-to-spherical maps
elsewhere, and by the determinant highest vector at infinity.

The global binary map is consequently a single scalar $S_i(m)$ times
these fixed intertwiners.  Include in $S_i(m)$ the conversion to CM
evaluation at a unit-framed object, namely
\[
 \prod_{\sigma:F\hookrightarrow\R}
       (2\pi i/\Omega_\sigma)^{2m},
\]
with the fixed algebraic unit powers prescribed in
Section~\ref{sec:theta}.  This scalar has bounded tame group-ring
coefficients.  Indeed, choose integral open tests at $q$ and fixed-level
open tests at $p$ that map a fixed compact vector to itself under
Equation~\ref{eq:binary-intertwiner}; the volume loss at $p$ is fixed,
whereas the full source maximal compact at $q$ has volume one.  The
binary moment estimate bounds the reference lift.  Its values are
$S_i(m)$ times values of one fixed nonzero algebraic compact automorphic
function, times determinant-character units.  Evaluate at a fixed
nonzero value, choosing finite representatives integral at $p$.
Division by that value costs a fixed denominator.  The cancellation
of determinant twists at finite compact levels prevents the source
summation level from growing with $q$.

Applying these intertwiners to the projected test at any boundary
component expresses its constant coefficient as $S_i(m)$ times a
compact automorphic evaluation and local operators.  Thus two estimates
remain: $S_i(m)$ must tend to zero through tame degree $b_*-1$, and the
projected local operators must have uniform bounds.  The latter will
use the same local factors as the period comparison.

\subsection{Raw toric periods and the comparison identity}

Let $f_0$ and $f_1$ be the weight-two newforms attached to $A_0$ and
$A_0^{D'}$, respectively, and let $\pi_0,\pi_1$ denote their unitary
automorphic representations.  For $m\ge1$, form toric periods over $K$ of
the $(m-1)$st normalized Maass--Shimura raises of these forms, against
$\eta_K^{-1}=(\cC^m\nu)^{-1}$.  Normalize CM evaluations by integral
differentials.  We denote the resulting periods by $P_{0,i}(m)$ and
$P_{1,i}(m)$, suppressing the tautological tame character from the
notation.  The following local choices fix their normalization.

\begin{enumerate}
\item At $q$, translate the spherical Whittaker vector by $n(1/q)$ in
torus coordinates whose first direction is the distinguished
multiplicative direction.  Multiply the fixed-measure torus integral by
$q-1$.  Thus $P_{j,i}$ is a raw sum at conductor $q$, with no averaging
denominator depending on $q$.
\item At $p$, use the Kirillov function
$W(\diag(x,1))=\1_{\Z_p^\times}(x)$.  This is the canonical orientation
of the $p$-depletion of the newvector.  The function belongs to the
Kirillov model of every generic local representation of $\GL_2$, including
ramified principal series, special representations, and supercuspidal
representations.
\item At all other places use fixed tests with nonzero local toric
pairing at weight zero.  For $P_{0,i}$ these can be the oriented Heegner
newvectors.  At a split level prime their diagonal Mellin integrals are
the nonzero local central $L$-factors.  For $P_{1,i}$, a prime dividing
$D'$ that is nonsplit in $K$ presents no obstruction: the local
representation is a tempered principal series, and the field torus acts
transitively on the projective line.  In the induced model this gives
the toric functional directly, since the central character is trivial.
A fixed deeper level is allowed at these finitely many places.
\end{enumerate}

The coefficients of both $P_{j,i}$ have a common lower valuation bound.
Indeed, all CM objects are ordinary with canonical multiplicative
structure, the $p$-tests are integral of fixed level, and Katz
differentiation gives bounded moments in integral differential frames.
The expansion and coefficient-multiplier results are
\citet[Section~5.2]{Katz1976} and
\citet[Theorem~6.10, Sections~11.1--11.7 and~15.21--15.24]{AndreattaGoren2005}.
The ordinary CM comparison is
\citet[Section~1.4, Theorem~1.2]{HidaTilouine1993}; its application
to the algebraized theta moments is Lemma~\ref{lem:unary-moments}.
Fixed CM-conductor and ideal-class denominators can be cleared once.
Changes to Igusa differential frames contribute fixed unit powers, and
the raw conductor-$q$ sum introduces no further denominator.  All these
assertions hold before projecting the tame group algebra to characters.

The relation between these periods and the binary scalar is the
following cleared identity.  Its augmentation condition is exactly
what will allow division after a fixed Artin truncation.

\begin{lemma}[Cleared period comparison]
\label{lem:cleared-period-comparison}
There are group-ring polynomial factors $A_i(m),B_i(m)$ with uniformly
bounded coefficients such that
\begin{equation}\label{eq:cleared-period-comparison}
 A_i(m)S_i(m)^2
   =B_i(m)P_{0,i}(m)^2P_{1,i}(m)^2.
\end{equation}
For each $i$, after taking $m$ sufficiently close to zero on its branch,
$A_i(m)$ has nonzero augmentation whose valuation is bounded above by
a constant independent of $i$.  The required neighborhood of zero may
depend on $i$.
The identity holds in the full characteristic-zero tame group algebra,
including at zeros of the global central $L$-values.
\end{lemma}

We prove this comparison below by matching the two bilinear period
formulas and their local factors.  First we show why $P_{0,i}$ is small.
Its square will then be small through twice the required tame order;
the bounded coefficients of $S_i$ will recover the required order for
$S_i$ itself.

\subsection{The Heegner logarithm limit}

We claim that
\begin{equation}\label{eq:toric-period-small}
 P_{0,i}(m)\longrightarrow0\pmod {t^{b_*}},
\end{equation}
where for each fixed $i$ one first takes $m$ sufficiently close to zero
on its branch, and then lets $i$ increase.

The passage from raised CM values to a Heegner logarithm follows the
$p$-adic period method of
\citet[Introduction and Proposition~3.24]{BertoliniDarmonPrasanna2013}.
Their construction is made under $p\nmid N$.  We give the primitive
and logarithm comparison for the reduction types occurring here.

Write $u=q_{\mathrm{Tate}}$ for a modular cusp parameter, to distinguish
it from the auxiliary prime $q_i$, and put
\[
 f_0^{[p]}(u)=\sum_{p\nmid n}a_nu^n,\qquad
 F_0(u)=\theta^{-1}f_0^{[p]}(u)
       =\sum_{p\nmid n}\frac{a_n}{n}u^n.
\]
On the ordinary Igusa tower, the normalized derivatives
$\theta^{m-1}f_0^{[p]}$ converge as $m\to0$ on the chosen branch, because
$n^{m-1}\to n^{-1}$ uniformly on the unit indices.  Their limiting
Igusa weight is zero.  The ordinary expansion principle therefore
constructs a global analytic function $F_0$ on the completed ordinary
locus of the good tame modular curve; weight-zero invariance descends
it from the Igusa tower.  Its differential and its cusp constants are
\begin{equation}\label{eq:inverse-katz-primitive}
 dF_0=f_0^{[p]}(u)\,d\log u,
 \qquad F_0|_{\text{each compatible cusp}}=0.
\end{equation}
Each component in use has such a compatible cusp.  The canonical
$\Gamma_0(p^{\ord_pN})$ section and the canonical-subgroup Frobenius
maps allow the $p$-depleted form to be read on this tame ordinary locus,
even when $p\mid N$.

To make the logarithm comparison explicit, choose a modular
parametrization $\varphi$ with
$\varphi^*\omega_0=\kappa f_0(u)d\log u$, where $\omega_0$ is a fixed
invariant differential of $A_0$ and $\kappa\ne0$ is fixed.  If the
depletion polynomial is $E_p(X)=\sum_r e_rX^r$, and $\mathrm{Fr}$ is
canonical-subgroup Frobenius, then
\begin{equation}\label{eq:depletion-differential}
 dF_0=\kappa^{-1}\sum_r e_rp^{-r}
       (\varphi\circ\mathrm{Fr}^{r})^*\omega_0.
\end{equation}
The factors $p^{-r}$ arise because substitution $u\mapsto u^{p^r}$
multiplies $d\log u$ by $p^r$.  There are only finitely many such maps,
all determined by the fixed newform.  They extend over a wide-open
thickening of the ordinary affinoid.

The primitive itself extends to a slightly smaller such thickening.
For completeness, clear one denominator and write its expansion near a
missing special-fiber point, with local parameter $z$, in the completed
localization of $\cO[[z]][1/z]$:
\[
 F_0=\sum_{n\in\Z}a_nz^n.
\]
Overconvergence of $dF_0$ implies exponential decay of the negative
coefficients $na_n$.  Division by $n$ loses only the subexponential
factor $|1/n|_p\le |n|$, so $a_n$ still decays exponentially on any
slightly smaller collar.  Exactness rules out a logarithmic term.  The
positive part is bounded and converges for $|z|<1$.  These expansions
retain the coefficients of the original $F_0$, hence introduce no
independent constants.  One can glue them on a model with collar
$zw=\varpi'$ and $|\varpi'|<1$ sufficiently close to one: modulo each
uniformizer power the usual localization--completion sequence
\[
 A\longrightarrow A[1/z]\oplus\widehat A
 \longrightarrow\widehat A[1/z]
\]
is injective at the first term and exact at the middle term.  This
follows from flatness of completion and invariance of local cohomology
with support.  Thus the continued primitive is an analytic function on
a connected wide open in every component under consideration.

We use Berkovich--Coleman integration on this neighborhood, with the
functoriality and fundamental theorem of
\cite[Definition~3.5]{KatzRabinoffZureickBrown2016}.  After a finite
extension, $A_0$ has either good or split multiplicative reduction.  In
the good reduction case integration of its invariant differential
equals its abelian logarithm
\cite[Corollary~3.18]{KatzRabinoffZureickBrown2016}.  In the
multiplicative case write $A_0^{\mathrm{an}}=\Gm^{\mathrm{an}}/Q_0^{\Z}$
and normalize the pullback differential to $dz/z$.  Starting with any
branch $\operatorname{Log}$, choose
\[
 \operatorname{Log}_{Q_0}(z)
 =\operatorname{Log}(z)
  -\frac{v_p(z)}{v_p(Q_0)}\operatorname{Log}(Q_0).
\]
This branch kills $Q_0$, and integration along every lifted path is the
abelian logarithm on the quotient; see
\cite[Example~3.11]{KatzRabinoffZureickBrown2016}.  All maps in
Equation~\ref{eq:depletion-differential} have the same elliptic target,
so this single branch suffices.

For a path from a compatible cusp $c$ to $x$, functoriality and
Equations~\ref{eq:inverse-katz-primitive}--\ref{eq:depletion-differential}
now give
\begin{equation}\label{eq:katz-abelian-log}
 F_0(x)=\kappa^{-1}\sum_r e_rp^{-r}
 \log_{A_0,p}\bigl(\varphi(\mathrm{Fr}^{r}x)\bigr).
\end{equation}
Indeed the cusp constant is zero, and every image of a cusp is torsion
by the Manin--Drinfeld theorem \citep{Manin1972,Drinfeld1973}.
This is an equality obtained by path integration on the connected wide
open; equality of derivatives alone would not suffice to remove
constants on different residue discs.

At the CM points entering $P_{0,i}$, the canonical Frobenius quotients
are the horizontal translations at $p$ of the conductor-$q_i$ Heegner
points, including when a nontrivial canonical $p$-level is present.
Taking their raw tame weighted sums in Equation~\ref{eq:katz-abelian-log}
therefore gives a fixed finite linear combination of the logarithm sums
in Proposition~\ref{prop:heegner-log}.  Both signs of the tame character
are covered by Equation~\ref{eq:heegner-log}.  The CM differential and
character unit powers tend to one as $m\to0$.  This proves
Equation~\ref{eq:toric-period-small}, with bounded higher coefficients
in every fixed truncation.  No ordinary hypothesis on $A_0$ at $p$ has
entered this argument.

\subsection{Period identities and the two local functional equations}

\begin{proof}[Proof of Lemma~\ref{lem:cleared-period-comparison}]
The equal-rank Rallis inner product formula and the Waldspurger toric
period formula give
\begin{equation}\label{eq:binary-period-comparison}
 S_i(m)^2=C_i(m)\,P_{0,i}(m)^2P_{1,i}(m)^2.
\end{equation}
We use these formulas with bilinear dual data, so the squares are
algebraic squares, rather than absolute squares.  For the theta
identity, the precise Siegel--Weil input is the anisotropic boundary
case of \cite[Theorems~2.1(i) and~2.2(i)]{Yamana2011}, also stated in
\cite[Theorem~7.1]{GanQiuTakeda2014}.  In the notation of those results,
$m=n=2$, $r=0$, and $s_0=(m-n)/2=0$.  On the doubled group $U(2,2)$
the Siegel Eisenstein series is holomorphic at $s_0$ and
$E(0,\Phi)=2I(\Phi)$, where $I$ is the ordinary convergent theta
integral over the compact binary group with quotient volume one.
The conversion to our global measure convention contributes only
the fixed scalar permitted in \cref{sec:conventions}.
Restrict this identity to
the two copies of $U(A)$ and integrate against the source cusp forms.
The theta seesaw gives their bilinear theta pairing on one side;
unfolding the Eisenstein series gives the doubling zeta integral on
the other.  Its Euler factorization is the central standard $L$-value
times the normalized local integrals used below.  This proves the
equal-rank Rallis identity needed here directly from that boundary
Siegel--Weil identity.  The doubling and local factorization
conventions are those of \cite[Sections~11.3, 11.5--11.6]{GanQiuTakeda2014}.
The source groups are compact, the standard base-change parameter is
cuspidal and tempered, and its zero-space theta lift vanishes, since
that lift could only be one-dimensional.  The Waldspurger identities
\cite[Section~II.7 and Section~III.3, Proposition~7]{Waldspurger1985} apply to cuspidal
transfers with compatible
central character and opposite torus characters in the two slots.
Thus the identities involve central values, not residues or derivatives.

For the theta pairing, pull back the same data by $(h,g)\mapsto(h^c,g^c)$.
Field conjugation is anti-symplectic on the tensor product and identifies
the oscillator with its dual by Poisson summation, without a global
conductor factor.  Pair fixed compact vectors nontrivially; if
necessary two different vectors can be used, since the scalar is the
same.  On the toric side, translate the vector by a rational torus
normalizer inducing conjugation.  This turns the second period into
the first, up to a fixed character value from a base-point translation.
At $p$ and $q$ the normalizer exchanges the distinguished coordinates
with unit scales.

The global central $L$-values in the formulas match by quadratic base
change over $K$:
\begin{equation}\label{eq:rankin-factorization}
 L\bigl(\tfrac12,\pi_L\otimes\eta_L\bigr)
 =L\bigl(\tfrac12,\pi_{0,K}\otimes\eta_K\bigr)
  L\bigl(\tfrac12,\pi_{1,K}\otimes\eta_K\bigr).
\end{equation}
Here $\eta_L$ is induced by norm from $\eta_K$, and simultaneous inversion
of the character gives the same comparison.  One may retain raw local
integrals at every removed place and use partial $L$-values.  Fixed
adjoint, norm, and global zeta factors then contribute a single fixed
scalar.  We compare the remaining local factors explicitly.

\paragraph{The matrix functional equation.}
Let $E$ be a split local field, write $\chi=\eta_v^{-1}$ in the
distinguished multiplicative coordinate, and let $\psi_{\mathrm{add}}$
be the trace additive character.  Write $\pi$ for the bare elliptic
local representation, so $\pi\simeq\pi^\vee$ and its central character
is trivial.  For a matrix coefficient $c$ and $F(X)=\Phi({}^tX)$, put
\[
 Z_G(F,c,\chi,s)=\int_{\GL_2(E)}
 F(g)c(g)\chi(\det g)|\det g|^{s+1/2}\,d^\times g.
\]
At $s=1/2$ this is a scalar coefficient of
Equation~\ref{eq:binary-intertwiner}.  The Fourier transform uses the
pairing $\psi_{\mathrm{add}}(\tr(XY))$.  Field conjugation exchanges the
two dual Lagrangians and acts on the source by inverse transpose, so
the second integral uses the full matrix Fourier transform and
$c^\vee(g)=c(g^{-1})$.  With these identifications the functional
equation is
\begin{equation}\label{eq:matrix-functional-equation}
 Z_G(\widehat F,c^\vee,\chi^{-1},1-s)
  =\gamma(s,\pi\otimes\chi,\psi_{\mathrm{add}})
     Z_G(F,c,\chi,s).
\end{equation}
Apply the matrix zeta identity
\cite[Theorem~13.1(iii),(v)]{JacquetLanglands1970} to
$\sigma=\pi\otimes(\chi\circ\det)$: its determinant exponent is
$s+1/2$ and its trace Fourier transform uses the self-dual additive
measure.  The $L$- and epsilon factors are the same as in the
Whittaker functional equation of Section~2.18 of that reference.
In particular, normalizing the
first map to be one on the chosen vector gives the factor
\begin{equation}\label{eq:local-gamma-factor}
 \Gamma_v(\pi\otimes\chi)
 =\varepsilon\bigl(\tfrac12,\pi\otimes\chi,\psi_{\mathrm{add}}\bigr)
   \frac{L(\tfrac12,\pi^\vee\otimes\chi^{-1})}
        {L(\tfrac12,\pi\otimes\chi)}.
\end{equation}
There is only one application of the functional equation in this
bilinear pairing.

To relate this calculation to the Rallis local integral, integrate
first over invertible matrices in $M_2(E)$.  Additive measure becomes a
fixed multiple of $|\det h|^2d^\times h$, so the integral factors as the
pairing of the two maps in Equation~\ref{eq:binary-intertwiner}.
For self-dual additive measure the spherical factor is
$(1-Q^{-1})(1-Q^{-2})$, where $Q$ is the residue cardinality.  It
cancels exactly when the local integral is normalized by its spherical
value: the spherical intertwiner itself gives
$L(1/2,\pi^\vee\otimes\chi)$ with no such factor.  In particular this
normalization never leaves a denominator $1-Q^{-1}$ at the moving
prime.

\paragraph{The toric functional equation.}
Put $a(x)=\diag(x,1)$, let
$w=\left(\begin{smallmatrix}0&1\\-1&0\end{smallmatrix}\right)$, and write
\[
 M(W,\chi,s)=\int_{E^\times}W(a(x))\chi(x)|x|^{s-1/2}\,d^\times x,
 \qquad W^w(g)=W(gw).
\]
The corresponding Whittaker functional equation is
\begin{equation}\label{eq:toric-functional-equation}
 M(W^w,\chi^{-1},1-s)
  =\gamma(s,\pi\otimes\chi,\psi_{\mathrm{add}})
    M(W,\chi,s).
\end{equation}
For the twist $\sigma$, use $W_\sigma(g)=\chi(\det g)W(g)$
and its central character $\chi^2$ in
\citet[Theorem~2.18(ii)--(iv)]{JacquetLanglands1970}.
Since $\det w=1$, the dual integral has character
$\chi(x)\chi(x)^{-2}=\chi(x)^{-1}$.
The identity $w a(x^{-1})=x^{-1}a(x)w$ then gives the displayed
Weyl translate, with no extra $\chi(-1)$ factor.  Both first integrals,
in Equations~\ref{eq:matrix-functional-equation}
and~\ref{eq:toric-functional-equation}, use the same $\chi$ in the same
distinguished coordinate.  The invariant Kirillov pairing factors the
toric local matrix integral into these Mellin integrals; its
proportionality constant also cancels under spherical normalization.
Consequently its factor is
\[
 \Gamma_v(\pi\otimes\chi)\,M(W,\chi,1/2)^2,
\]
with exactly the same gamma orientation as the matrix factor.

At $p$, the character $\chi$ is unramified and the chosen Kirillov test
gives $M(W,\chi,1/2)=1$.  Thus all gamma, Euler, and epsilon factors
cancel between the two period formulas.  No bound for
$\Gamma_v(\pi\otimes\chi)$ itself is asserted or needed; indeed
$v_p(\chi(p))=m$ varies with the weight.  The functional equations are
identities of rational zeta functions and apply by continuation at
these parameters.  Their validity includes special representations
\cite[Proposition~3.6 and Corollary~3.7]{JacquetLanglands1970}; a
fixed compact open test normalizes the matrix map even when $\pi_p$ is
ramified.  Additive-character signs and changes of the chosen unit
frames introduce only character values on units and fixed unit
powers.  Since $p$ splits completely in $L$, there are two matching
local comparisons; the quadratic twist by $D'$ does not change either
local elliptic representation.

\paragraph{The raw factor at $q$.}
Let $L_{\mathrm{fake}}$ be the degree-two unramified central Euler
factor formed using $\chi(q)$ while ignoring its tame inertia.  The
spherical Whittaker expansion after translation by $n(1/q)$ gives
\begin{equation}\label{eq:raw-q-factor}
 F_q(\chi):=(q-1)M(W,\chi,1/2)
 =\sum_{a\in\F_q^\times}\chi(a)
       \bigl(\psi_{\mathrm{add}}(a/q)+L_{\mathrm{fake}}-1\bigr).
\end{equation}
To obtain the formula, separate valuation zero from positive
valuations.  On units the translated Whittaker function has the
additive factor $\psi_{\mathrm{add}}(a/q)$; on positive valuations its
sum is $L_{\mathrm{fake}}-1$.  Multiplication by $q-1$ converts the
unit integrals to the displayed raw sum.  For nontrivial tame $\chi$
the second sum vanishes, leaving the Gauss sum.  At augmentation one
instead has
\[
 F_q(1)=-1+(q-1)(L_{\mathrm{fake}}-1).
\]
In the group ring put $D_q=L_{\mathrm{fake}}^{-1}$.  Clearing this
Euler denominator gives the polynomial factor
\begin{equation}\label{eq:cleared-q-factor}
 H_q=D_qF_q
   =D_q\sum_{a\in\F_q^\times}\chi(a)\psi_{\mathrm{add}}(a/q)
     +(1-D_q)\sum_{a\in\F_q^\times}\chi(a),
 \qquad H_q(1)=(q-1)-qD_q(1).
\end{equation}
Both $D_q$ and $H_q$ have bounded group-ring coefficients.  The
no-eigenvalue-one condition on $V(\gamma)$ implies that the limit of
$D_{q_i}(1)$ at weight zero is nonzero; the conventions of
Section~\ref{sec:conventions} identify this with the corresponding
Frobenius Euler determinant, up to factors tending to one.  Thus its
valuation is bounded above.  Since $q_i\to1$, the same is true of
$H_{q_i}(1)$, once $m$ is sufficiently close to zero for each $i$.
The two raw toric periods contribute $F_q^2$ in their respective
comparisons.  Therefore the ratios in
Equation~\ref{eq:binary-period-comparison} divide by these full factors,
and never by an isolated $(q-1)^2$, by a conductor projector, or by
individual Gauss sums.  The identical gamma factors cancel just as at
$p$.

\paragraph{The archimedean factors.}
Write $k_0=m-1$.  In the $2\times2$ Fock model the Gaussian norm ratio
of $\det(Z)^{k_0}$ to the vacuum is
\begin{equation}\label{eq:factorial-cancellation}
 k_0!(k_0+1)!.
\end{equation}
For example, expand the determinant by the binomial theorem.  The
different monomials are orthogonal, and each of the $k_0+1$ summands
contributes $(k_0!)^2$ after including its squared binomial coefficient.
Their sum is Equation~\ref{eq:factorial-cancellation}.  The definite
quaternion real polarization has positive quadratic norm equal to this
determinant.  Its scalar symplectic $\SL_2$ generates from the vacuum
the holomorphic discrete series of lowest weight two, and raising is
multiplication by the determinant in Fock coordinates.  The norm ratio
of its $k_0$th raised vector is the same $k_0!(k_0+1)!$.  This is also
the ratio in the archimedean toric pairing with the Weyl-dual raised
weight-two vector.

The factorials therefore cancel.  The remaining comparison with
normalized Maass--Shimura raising is a fixed algebraic $p$-unit to the
$k_0$th power: this follows by comparing the leading linear
Schr\"odinger coordinates, as in Lemma~\ref{lem:unary-moments}, with
quaternion frames chosen as unit bases at $p$ and CM heights chosen as
algebraic $p$-units.
The anti-symplectic duality contributes only phases of the same kind.
There is one comparison at each real embedding of $F$, and the
compatible CM normalization on both sides has weight $2m$ per
embedding, with the factors $(2\pi i/\Omega_\sigma)^{2m}$ already
included in $S_i(m)$.  No factorial denominator remains.

\paragraph{Fixed factors and zeros.}
At fixed exceptional finite places away from $p$, the local factors
are rational functions of unramified smooth-character parameters,
with fixed inertia.  On our branch those parameters tend to their
weight-zero values.  At the tempered center the zeta denominators have
no pole, and the toric functionals were chosen nonzero there.  After
shrinking the branch, the denominators have nonzero specialization of
bounded valuation, so these ratios cost a fixed denominator.
The same statement applies to the fixed rational local intertwiners.
The remaining global measure and norm constant is independent of
$i,m,\nu$.  It is algebraic by evaluating the comparison at any
parameter for which both sides are nonzero, since all normalized
periods there are algebraic.  If there is no such parameter, the
vanishing theta comparison already gives the required zero scalar.

The factors just computed, notably $D_q$ and $H_q$, supply the
polynomial clearing factors in
Equation~\ref{eq:cleared-period-comparison}, with bounded coefficients
and with nonzero augmentation of $A_i(m)$ of bounded valuation.
Fixed algebraic constants and unit powers are
absorbed into $A_i$ and $B_i$.  The period formulas prove this identity
first at every finite-order character.  Clearing denominators gives
an identity in the full characteristic-zero tame group algebra,
because those characters separate that algebra.  There is no division
by a global central $L$-value: use
Equation~\ref{eq:rankin-factorization} in the two identities and
cross-multiply the local factors.  Hence the cleared identity also
holds at their global zeros.  We do not assert that the inverse of
$A_i$ is a bounded measure on the entire finite group.

\end{proof}

\subsection{Projected tests at every cusp}

The period comparison now controls the scalar shared by all binary
theta values.  To apply it to every constant coefficient, we must bound
the projected local operators in the factorization at the start of
this section, uniformly in the boundary representative.  The moving
prime will use the factor $D_q$ just computed.

At $p$, canonical projection leaves the binary test supported on
integral matrices.  Expanding Equation~\ref{eq:binary-intertwiner} by
the determinant valuation gives the factor
\[
 p^{(m-1)v_p(\det h)}.
\]
The right-coset volumes have fixed level, and translates of the fixed
compact automorphic function have bounded sup norm.  For large $m$
the integral-matrix expansion therefore converges $p$-adically with a
uniform bound.  Its determinant truncations are the usual
Godement--Jacquet series, so by continuation it computes the same
rational zeta operator.  This argument applies at both places over
$p$, without imposing any condition on the local elliptic
representation.

At $q$, the projected test is supported integrally, depends only on
reduction modulo $q$, and retains determinant equivariance under the
source $\GL_2(\cO_{F_v})$.  On the full-rank stratum this gives a bound
directly with maximal-compact volume one.  On a deficient-rank stratum
the determinant of the stabilizer surjects onto $\Delta_i$.  Therefore
its group-ring coefficient is invariant under multiplication by every
element of $\Delta_i$ and is a bounded multiple of the norm element
\[
 N_{\Delta_i}=\sum_{\delta\in\Delta_i}[\delta].
\]
This is an integral invariance statement: all its group coefficients
are equal.  It involves no division by $|\Delta_i|$.

The remaining integral, now untwisted on units, is a sum over integral
image lattices $h\cO_{F_v}^{2}$ with weight
$|\det h|\chi(\det h)$, and a coefficient constant on each prescribed
reduced image subspace.  For a line or the zero subspace, express the
condition of having that exact reduced image by subtraction of the
conditions that the image be contained in it.  Each containment sum is
a translate and a determinant power times the spherical full-lattice
series $L_{\mathrm{fake}}$.  All translation determinants and volume
factors are powers of $q$, hence $p$-units.  Thus the already controlled
factor $D_q$ clears the operator denominator.  This rank-stratum
argument is valid coefficientwise in the full group ring.

At all other places the operators are spherical or among fixed
rational operators with the no-pole property proved above.  Boundary
Iwasawa representatives can be chosen integral at $p$, compatible with
the canonical multiplicative flag.  Evaluation of the compact model
at the conjugate representatives then introduces only unit factors.
Source sums have a fixed volume denominator, and all nonsplit
spherical lines and fixed-level coefficients have common bounds.
Although the number of boundary representatives may grow at $q_i$,
their projected tests have precisely the uniformly bounded forms just
described.  This proves uniformity over every cusp, rather than only
over a chosen finite list of evaluations.

\subsection{Truncation, square roots, and choice of weights}

We finish by spelling out the algebra behind the coefficientwise
estimate.  Fix $h\ge2b_*$.  In a coefficient extension containing the
finite test values, substitute
\[
 [g]\longmapsto(1+t)^{\ell_i(g)}\pmod {t^h}
\]
in the tame group algebra.  At finite precision this is a ring map
modulo $p^{n_i-c_h}$ for a constant $c_h$ depending only on $h$:
the relation $[g]^{p^{n_i}}=1$ maps to
$(1+t)^{p^{n_i}}-1$, whose coefficients below degree $h$ have valuations
at least $n_i-c_h$.  Consequently products and identities may be
truncated with an error tending to zero.  Bounded group-ring
coefficients give bounded coefficients in this fixed Artin
truncation.  All estimates below are unaffected by these errors.

If $D(t)=d_0+\cdots+d_{h-1}t^{h-1}$ has bounded coefficients and
$v_p(d_0)$ is bounded above, its inverse in the truncated rational ring
is bounded by a constant depending only on $h$ and those bounds:
\begin{equation}\label{eq:bounded-artin-inverse}
 D(t)^{-1}=d_0^{-1}
  \sum_{r=0}^{h-1}\bigl(-d_0^{-1}(D(t)-d_0)\bigr)^r
  \pmod {t^h}.
\end{equation}
Thus the clearing factors must be cleared in the group algebra
\emph{before} truncation, after which their inverses cost a fixed
denominator.  No assertion about bounded character idempotents or
bounded inverses on the full group ring is required.

By Equation~\ref{eq:toric-period-small}, the coefficients of $P_{0,i}$
below degree $b_*$ tend to zero, and all coefficients below degree $h$
of both toric periods are bounded.  Every coefficient of $P_{0,i}^2$
below degree $2b_*$ contains a factor of degree below $b_*$.  Hence
\[
 P_{0,i}^2P_{1,i}^2\longrightarrow0\pmod {t^{2b_*}}.
\]
Apply Equation~\ref{eq:bounded-artin-inverse} to
Equation~\ref{eq:cleared-period-comparison} to obtain
$S_i(m)^2\to0\pmod {t^{2b_*}}$.

The passage to the square root uses only valuations, not compactness
of the fields of coefficients.  Write $S_i=\sum_{j<h}s_{ij}t^j$, with
all $s_{ij}$ bounded.  First $s_{i0}^2\to0$, hence $s_{i0}\to0$.
Inductively, for $0<j<b_*$,
\[
 [t^{2j}]S_i^2
  =s_{ij}^2+2\sum_{a<j}s_{ia}s_{i,2j-a}.
\]
The sum tends to zero by the preceding steps and boundedness, so
$s_{ij}^2\to0$ and then $s_{ij}\to0$.  Therefore
$S_i(m)\to0\pmod {t^{b_*}}$.  The same fixed-truncation inversion
handles the projected-test factors at every cusp.  Their uniform
bounds prove the asserted vanishing of all constant coefficients.

Finally, for each $i$, let $m$ approach zero in the prescribed branch
until the interpolation of $P_{0,i}$ has the precision supplied by
Equation~\ref{eq:heegner-log}, and until all fixed and moving clearing
factors satisfy the bounds above.  Arbitrarily large positive integers
exist in every sufficiently small such congruence class.  Choose one
that also satisfies any imposed archimedean lower bound and additional
$p$-adic approximation requirements.  Taking the precision to infinity
with $i$ gives the required diagonal choice of weights, and proves
Proposition~\ref{prop:constant-term-small}.
